\documentclass[../../main.tex]{subfiles}
\begin{document}

\begin{center}
  {\fontsize{48pt}{58pt}\selectfont \textbf{1968}\normalsize\textbf{年度}}
\end{center}

\begin{table}[htb]
  \centering
  \caption{出題テーマと難易度}
  \begin{tabular}{|c|p{9cm}|c|}
    \hline
    \textbf{問題番号} & \textbf{テーマ} & \textbf{難易度} \\
    \hline
    第1問 & 二曲線の間の領域を二等分する直線 & \\
    \hline
    第2問 & 正四角錐の斜面がなす二面角 & \\
    \hline
    第3問 & 二次関数の係数と二点の値の組 & \\
    \hline
    第4問 & 回転曲面で囲まれる立体の体積 & \\
    \hline
    第5問 & 条件を満たす整数係数の三次多項式 & \\
    \hline
    第6問 & 分数冪の不等式と大きな整数の桁数 & \\
    \hline
  \end{tabular}
\end{table}

\begin{table}[htb]
  \centering
  \caption{解答}
  \renewcommand{\arraystretch}{1.4}
  \begin{tabular}{|c|c|p{9cm}|}
    \hline
    \textbf{問題番号} & \textbf{小問} & \textbf{答え} \\
    \hline
    第1問 & -- & $\alpha=-1$ \\
    \hline
    第2問 & -- & 二面角は$\dfrac{\pi}{3}$ \\
    \hline
    第3問 & -- & $\alpha=\beta=1$なら原点，$\alpha=1\ne\beta$なら$y$軸，$\beta=1\ne\alpha$なら$x$軸，$\alpha=\beta\ne1$なら直線$y=x$，それ以外なら直線$y=\dfrac{\beta-1}{\alpha-1}x$を除いた平面 \\
    \hline
    第4問 & -- & $V=\dfrac{46\pi}{15}$ \\
    \hline
    第5問 & -- & $P(x)=4x^3+x$ \\
    \hline
    第6問 & (i) & $1+\dfrac{\alpha x}{2}\le(1+x)^\alpha$（証明問題） \\
    \cline{2-3}
    & (ii) & $2^{200}$は$61$桁，最上位は$1$ \\
    \cline{2-3}
    & (iii) & $0.300<\log_{10}2<0.302$（証明問題） \\
    \hline
  \end{tabular}
\end{table}

\paragraph{試験形態}
6問．試験時間は要確認．

\end{document}
